% TOP_PROBLEM: {"id": 3266, "title": "Long directed cycles in diregular digraphs", "category_id": 3, "category": {"id": 3, "name": "graph_theory", "display_name": "Graph Theory", "description": "Problems involving graphs, networks, and their properties.", "slug": "graph-theory", "order_index": 3, "created_at": "2026-07-31T15:26:25.671Z"}, "status": "open", "problem_number": "OPG-46460", "set_id": 12, "statusQualification": "The mathematical statement uses minimum in- and outdegree, not exact regularity suggested by the title."}
% FIRST_POSTED: 2026-10-02
% ATTEMPT_STATUS: unresolved
% UnsolvedMath ID 3266; source catalogue code OPG-46460
% !TeX program = lualatex
% GPT-6 Astra Ultra research attempt; independently checked partial work, 2026-10-02.
\documentclass[11pt,a4paper]{article}
\usepackage[margin=25mm]{geometry}
\usepackage{fontspec}
\setmainfont{lmroman10-regular.otf}[BoldFont=lmroman10-bold.otf,
 ItalicFont=lmroman10-italic.otf,BoldItalicFont=lmroman10-bolditalic.otf]
\usepackage{amsmath,amssymb,mathtools}
\usepackage{unicode-math}
\setmathfont{latinmodern-math.otf}
\AtBeginDocument{\let\setminus\smallsetminus}
\usepackage{xurl,hyperref}
\hypersetup{colorlinks=true,urlcolor=blue}
\setcounter{secnumdepth}{0}
\setlength{\parindent}{0pt}
\setlength{\parskip}{5pt}
\setlength{\emergencystretch}{3em}
\begin{document}
\section*{Long directed cycles in diregular digraphs}\label{3266}
{\small UnsolvedMath ID 3266 · OPG-46460 · Graph Theory · OpenGarden}
\subsection{Definitions and mathematical statement}
Let $D=(V,A)$ be a finite oriented graph: it has no loops, no parallel
arcs, and at most one of $uv,vu$ for each pair of distinct vertices.
The indegree $d^-(v)$ counts arcs entering $v$, and the outdegree $d^+(v)$
counts arcs leaving it. The digraph is strongly connected if for every
ordered pair $u,v$ there is a directed path from $u$ to $v$.

A directed cycle of length $\ell$ consists of $\ell$ distinct vertices
$v_1,\ldots,v_\ell$ and the arcs $v_iv_{i+1}$, with indices read modulo
$\ell$. The conjecture asks whether, for every integer $d\ge1$, every
strongly connected oriented graph satisfying
\[
                  d^+(v)\ge d\quad\text{and}\quad d^-(v)\ge d
                  \qquad(v\in V)
\]
contains a directed cycle of length at least $2d+1$.

These are minimum-degree hypotheses. Despite ``diregular'' in the
catalogue title, the displayed source statement does not require all
indegrees and outdegrees to be equal; that narrower interpretation
would lose part of the target. There is also no upper bound on the
number of vertices, so the requested cycle need not be Hamiltonian.

The orientation assumption implies $|V|\ge2d+1$, because a vertex has
disjoint in- and outneighbour sets. Thus the requested length is compatible
with the necessary order bound. The strong-connectivity assumption is
retained explicitly: it supplies directed paths between all parts of
the graph, not merely connectedness after forgetting directions.
\subsection{Short English statement}
Does strong connectivity together with minimum indegree and outdegree d force an oriented graph to contain a directed cycle on at least 2d+1 vertices?
\subsection{Research attempt: the tournament case and a sharp strong example}
\paragraph{Scope and literature check.}
GPT-6 Astra Ultra, 2 October 2026. Both degrees are bounded below;
they need not be equal. Jackson's paper [S3] is primary background
for this long-cycle question. Lo--Patel--Y\i ld\i z [S4] prove a
Hamiltonicity result for sufficiently large regular oriented graphs
in an order range including $n\le4d+1$. Their exact regularity
assumption is stronger than the minimum-degree assumptions here.
The following argument proves the tournament subclass and constructs
strong, nonregular examples attaining precisely the proposed length.

\paragraph{1. An elementary bound using one endpoint.}
Let $v_0,\ldots,v_\ell$ be a longest directed path. All at least $d$
outneighbours of $v_\ell$ lie on the path, and none is $v_{\ell-1}$,
since opposite arcs are forbidden. The earliest has index at most
$\ell-d-1$. Closing the corresponding terminal segment produces a
directed cycle of length at least $d+2$. In particular the full
$2d+1$ target follows when $d=1$. This elementary bound neither uses
strong connectivity nor minimum indegree, and that unused information
identifies a natural place where a stronger argument must enter.

\paragraph{2. Strong tournaments have a spanning directed cycle.}
For completeness, choose a longest directed cycle $C$ in a strong
tournament. Such a cycle exists because every vertex can reach every
other vertex and the graph is nontrivial. Suppose some vertex $x$
outside $C$ has both an incoming and an outgoing adjacency with $C$.
Around the cyclic order there is then a transition
$c_i\to x\to c_{i+1}$. Inserting $x$ gives a longer cycle, impossible.
Thus every outside vertex either dominates all of $C$, forming a
set $A$, or is dominated by all of $C$, forming a set $B$.

If outside vertices exist, strong connectivity implies both $A$ and
$B$ are nonempty: with only $A$, no arc leaves $C$, and with only
$B$, no arc enters $C$. A directed path from $C$ to $A$ must first
leave $C$ into $B$ and later cross from $B$ to $A$. In particular
some arc $b\to a$ exists with $b\in B,a\in A$. Replacing any arc
$c_i\to c_{i+1}$ of $C$ by
$c_i\to b\to a\to c_{i+1}$ again gives a longer cycle.
The contradiction proves that $C$ is Hamiltonian.

A tournament satisfying the degree hypotheses has at least $2d+1$
vertices, because the in- and outneighbour sets at each vertex are
disjoint. The Hamilton cycle therefore proves the required conclusion
for every strong tournament. Missing adjacencies are exactly what
prevents the two-set classification of outside vertices in an
arbitrary oriented graph.

\paragraph{3. Sharpness under the actual strong-connectivity hypothesis.}
For any $d\ge1$, take two copies of the cyclic regular tournament
on $2d+1$ vertices and identify exactly one vertex, called $z$.
Within each copy the arcs are $x\to x+s$ modulo $2d+1$ for
$1\le s\le d$. Put no arcs between the other vertices of different
copies. The result is an oriented graph on $4d+1$ vertices.
Nonidentified vertices have indegree and outdegree $d$, while $z$
has both degrees $2d$. Unit-step Hamilton cycles show that each
copy is strong, and paths through $z$ show that their union is strong.

A simple directed cycle cannot use nonidentified vertices from both
copies. Traversing from one copy to the other and returning would
force visiting $z$ twice. Every cycle therefore lies in one copy,
so has length at most $2d+1$. Each copy supplies a cycle of that
length. This is an exact equality family satisfying all the
canonical hypotheses, including strong connectivity. It also shows
why a theorem requiring every degree to equal $d$ cannot simply be
applied to the minimum-degree version.

\paragraph{Verification and first remaining gap.}
The script [S9] constructs these examples for $d=1,2,3$, verifies
both degree bounds and strong connectivity, and enumerates simple
directed cycles to confirm their maximum lengths. The proofs above
hold for every $d$. Beyond tournaments, an outside vertex may have
missing adjacencies to the cycle; neither insertion nor the uniform
$A,B$ partition follows. A mechanism exploiting the two minimum-degree
bounds despite those missing pairs is still absent.

\paragraph{Final status.}
UNRESOLVED for general strongly connected oriented graphs. The
tournament subclass, the case $d=1$, and sharpness examples are proved.
\subsection{Sources}
\begin{itemize}
\item[S1] ULAM.AI and the UnsolvedMath Contributors, supplied problems.json, record 3266 (OPG-46460), OpenGarden collection. Catalogue identity and source transcription; CC BY 4.0.
\url{https://huggingface.co/datasets/ulamai/UnsolvedMath}.
\item[S2] Open Problem Garden, Long directed cycles in diregular digraphs. Original problem and discussion; the mathematical formulation below is expanded from the supplied source record.
\url{http://www.openproblemgarden.org/op/long_directed_cycles_in_digraph_with_minimum_in_and_out_degree}.
\item[S3] Bill Jackson, \emph{Long paths and cycles in oriented graphs}, Journal of Graph Theory 5 (1981), 145--157. \url{https://doi.org/10.1002/jgt.3190050204}.
\item[S4] Allan Lo, Viresh Patel and Mehmet Akif Y\i ld\i z, \emph{Cycle Partitions in Dense Regular Digraphs and Oriented Graphs}, Forum of Mathematics, Sigma (2025). \url{https://arxiv.org/abs/2309.11677}; \url{https://doi.org/10.1017/fms.2025.28}.
\item[S9] GPT-6 Astra Ultra, exact finite checks accompanying this attempt (2 October 2026). Repository script, Python standard library only. \texttt{scripts/check\_attempt\_batch03\_digraphs\_2026\_10\_02.py}.
\end{itemize}
\end{document}
